\section{Implementation specification}

\subsection{Minimal software interfaces}
The first implementation can remain framework-agnostic. The essential interfaces are:
\begin{lstlisting}[language=Python,basicstyle=\ttfamily\small,frame=single,breaklines=true]
class RouteMetric:
    def distance(self, r: Tensor, s: Tensor) -> Tensor: ...

class FunctionalDescriptor:
    def fit_batch(self, host_states: Tensor, synth: Tensor) -> Tensor: ...

class SPDGeometry:
    def distance(self, c1: Tensor, c2: Tensor) -> Tensor: ...
    def log_map(self, base: Tensor, target: Tensor) -> Tensor: ...
    def exp_map(self, base: Tensor, tangent: Tensor) -> Tensor: ...
    def parallel_transport(self, source: Tensor, target: Tensor,
                           tangent: Tensor) -> Tensor: ...

class MetricController:
    def parameters_regularizer(self) -> Tensor: ...
    def diagnostics(self) -> dict[str, float]: ...
\end{lstlisting}
Every implementation must expose conditioning and repair diagnostics; silent projection to SPD is prohibited in qualification runs.

\subsection{Numerical safeguards}
\begin{enumerate}[leftmargin=1.7em]
  \item symmetrize $C\leftarrow(C+C^\top)/2$ before decomposition;
  \item use a declared eigenvalue floor and report how often it is active;
  \item compare eigenvalue and Cholesky implementations on the same inputs;
  \item use double precision for operator verification, then benchmark mixed precision separately;
  \item check $\Exp_C(\Log_C(D))\approx D$ and transport norm behavior;
  \item monitor learned-map monotonicity and inverse round-trip error;
  \item reject batches with undefined route probabilities rather than masking them.
\end{enumerate}

\subsection{Proposed repository placement}
This article is integrated into the existing \texttt{geometry-mmalls-g1} repository under:
\begin{verbatim}
paper/riemannian_learning_architecture/
  main.tex
  references.bib
  README.md
  Makefile
  figures/
docs/reports/
  MMALS_Riemannian_Learning_Architecture_v1_0.pdf
docs/specs/
  MMALS_Riemannian_Experimental_Protocol_v1_0.md
governance/riemannian_learning_architecture/
  source_tiers.yaml
  claim_manifest.yaml
  uploaded_source_hashes.txt
scripts/
  validate_riemannian_article.py
releases/
  MMALS_Riemannian_Learning_Architecture_v1_0_full_package.zip
.github/workflows/
  riemannian-article.yml
\end{verbatim}


\section{Source tiers and claim governance}

\subsection{Tier definitions}
\begin{center}
\begin{tabularx}{\textwidth}{@{}llX@{}}
\toprule
Tier & Label & Meaning \\
\midrule
T1 & Established primary research & Peer-reviewed papers, standard mathematical monographs, or independently established results directly supporting a component. \\
T2 & Integrated primary synthesis & The uploaded 2026 doctoral thesis synthesizing and extending the underlying Riemannian deep-learning papers. \\
T3 & Executed MMALS evidence & Versioned notebooks, result bundles, reports, and repository evidence produced by the project. Valid only for their declared benchmark and protocol. \\
T4 & Proposed MMALS mechanism & Architecture, loss, hypothesis, or roadmap introduced here and not yet empirically qualified. \\
T5 & Speculative future branch & Hyperbolic hierarchy, density-state memory, or cross-domain extension lacking prerequisite evidence. \\
\bottomrule
\end{tabularx}
\end{center}

\subsection{Claim ladder}
\begin{center}
\begin{tabularx}{\textwidth}{@{}llX@{}}
\toprule
Code & Status & Minimum evidence \\
\midrule
C0 & Build validity & PDF compiles, schemas validate, numerical operator tests pass. \\
C1 & Diagnostic validity & Intrinsic metric adds out-of-sample explanatory value over Euclidean diagnostics. \\
C2 & Controlled utility & One geometry-aware component improves a declared endpoint on controlled factors at matched budget. \\
C3 & Continual utility & Improvement replicates with retention and plasticity jointly reported. \\
C4 & Cross-benchmark robustness & Effect repeats on natural-image or domain streams. \\
C5 & Operational geometry & Geometry changes routing/control decisions for the better under cost and calibration constraints. \\
C6 & General architectural claim & Multiple independent replications support a reusable \MMALS{} geometric mechanism. \\
\bottomrule
\end{tabularx}
\end{center}
The current article is C0 as a package and T4 as a mechanism proposal. It imports T1/T2 components and T3 evidence but creates no new executed model result.

